\chapter*{Abstract} % no number
\label{abtract}

\addcontentsline{toc}{chapter}{Abstract} % add to index

Photovoltaic power forecasting is usually posed as a per-plant supervised problem: one
model is fitted to the history of one installation. That formulation does not survive
contact with deployment, where a newly connected plant arrives with a couple of weeks of
measurements and no training set of its own. This thesis studies the alternative
formulation --- \emph{cross-plant zero-shot forecasting}, in which a model trained on one
fleet must forecast plants it has never observed --- and asks whether frozen foundation
models, fused across modalities, are the right instrument for it.

The work makes three contributions. First, it constructs a dataset of record for the
problem: 1.34 million rows over 110 photovoltaic sites in two regions, pairing measured
power and fourteen weather and solar-geometry covariates with co-registered satellite
imagery on a single timestamp grid, together with a plant-disjoint split and a tensor
contract that every model in the study consumes unchanged. Second, it defines a fairness
protocol --- identical physical-time windows, identical splits, identical
capacity-normalised metrics, no domain physics inside models --- and runs \emph{28}
forecasters under it, spanning statistical references, classical machine learning,
supervised deep networks, four families of time-series foundation model, retrieval and
adapter-based adaptation of frozen backbones, and both endogenous and exogenous multimodal
systems. Third, it proposes MMTSFM, an architecture that fuses a frozen Chronos-2
time-series backbone with a frozen V-JEPA~2.1 video encoder through two mechanisms:
\emph{selective temporal interleaving}, which weaves visual tokens into the time-series
sequence only inside a short refinement window at roughly two percent token overhead, and
\emph{causal Grassmann mixing}, an $O(L)$ temporal operator that encodes the transition
between consecutive hidden states as a subspace on a Grassmann manifold rather than as a
dot-product similarity.

The benchmark yields four findings that reframe the problem. A well-tuned supervised
transformer (iTransformer, skill score $0.552$) remains the model to beat, ahead of every
foundation model evaluated. Zero-shot time-series foundation models underdeliver on
photovoltaic dynamics, one of them falling below the naive reference. Retrieval over a
frozen backbone reaches skill $0.478$ against $0.331$ for fine-tuning the same backbone,
making it the most effective and cheapest adaptation tested. And the best multimodal model
in the suite synthesises its images \emph{from the series itself}, while every system
consuming genuine satellite frames sits mid-pack --- evidence that visual inductive bias,
not visual information, is what currently converts into forecast skill. MMTSFM is presented
in full, with its training curriculum, its controls and its evaluation plan; the
experimental campaign for the model was still running at the time of writing, and the
chapter that reports it states the measurement protocol and the placeholders rather than
provisional numbers.

\par

